\guidepart{Final Chapter}{
This part mainly includes "Change and Constants in Technology" and "Afterword."

"Change and Constants in Technology" invites us to stop chasing the next headline and return to a slower timescale, to discern the undercurrents that still govern the world beneath the churn of tools. This way, no matter which way the wind blows, we know what is worth pursuing and what must be held fast.

\vspace{1em}

The "Afterword" records the letters exchanged between my daughter and me during the writing process. She found this book interesting, un-pedantic, and rhythmic in its mix of long and short sentences. These "rainbow reviews" gave me tremendous strength and motivation to keep writing. In my replies to her, I also shared my thoughts on choosing schools and choosing majors. As a father, I wanted to use this book to tell her everything I consider important---but if I could say only one thing, what would it be?

\vspace{1em}

The closing section of the "Afterword" contains the acknowledgments for this book. This is a book that was pushed into existence. Without the pain of being torn apart, without the feedback my friends spent so much time writing for me, and without the help of every person named in the acknowledgments, this book would still be nothing more than an idea in my head.
}


\chapter*{Change and Constants in Technology}
\phantomsection
\addcontentsline{toc}{chapter}{Change and Constants in Technology}
\markboth{Change and Constants in Technology}{}
\thispagestyle{plain}

\begin{myquote}

If we observe from the perspective of change, then heaven and earth cannot endure even a single blink;

if we observe from the perspective of the unchanging, then all things and ourselves are equally infinite.

\hfill ——\textit{The First Ode to the Red Cliff}

\end{myquote}

\section{Three Threads of the Fifth Dimension}

Surveying the entire journey of the Fifth Dimension, we have witnessed an extraordinary amount of change.
Teleportation compressed distance, the Stealth Codex accelerated flow, the Absorption Skill rewrote the very definition of capability\ldots Technology races ahead, paradigms shift, and yesterday's kings become today's discards in the blink of an eye.

But if we stretch the timescale long enough---from the steam engine to the internet, from the abacus to quantum computing, from single-celled organisms to human civilization---we discover three threads running deep beneath history, governing the evolution of all complex systems like hidden currents.

\textbf{The first: the simple always defeats the complex.}

\textbf{The second: the flowing always defeats the hoarded.}

\textbf{The third: the self-evolving always defeats the static.}

\vspace{1em}

These three threads are not independent. They are three facets of a single truth:

Simplicity is the \textbf{precondition} for evolution---complex systems are too fragile to withstand the shock of mutation.

Flow is the \textbf{fuel} of evolution---closed systems can only spin in circles around local optima.

Evolution is the \textbf{ultimate form} of simplicity and flow---a system that is simple enough and open enough will spontaneously move toward self-improvement.

\vspace{1em}

Understanding these three threads helps us find \textbf{unchanging anchors} amid technological upheaval---

No matter which way the wind blows, we know where to stand.

\newpage
\section{The First Thread: Simplicity Defeats Complexity}

% \begin{myquote}
% Simplicity is the ultimate sophistication.

% \hfill ——达・芬奇
% \end{myquote}
%——————————————————————————————————————
\subsection*{The Bitter Lesson}
%——————————————————————————————————————
In 2019, reinforcement learning pioneer Rich Sutton wrote a short essay that sent shockwaves through the AI community for years to come. He called it \textbf{The Bitter Lesson} (\textit{The Bitter Lesson})\footnote{Richard Sutton, \textit{The Bitter Lesson}, 2019.}

\begin{myquote}
AI researchers have always tried to encode human knowledge into their systems.

In the short term, this always helps, and it gives researchers a sense of satisfaction.

But in the long run, it hinders progress.

The real breakthroughs come from the opposite direction---

scaling up computation through search and learning.
\end{myquote}

In other words: \textbf{skip the clever tricks; just stack more compute.}

\vspace{1em}

This conclusion was hard for many AI researchers to accept.

The feature engineering, knowledge graphs, and rule-based systems they had spent decades meticulously designing were swept away by brute force.

\vspace{1em}

In chess, Deep Blue relied on opening books and endgame tables distilled from human masters---elegant, refined, brimming with human wisdom. In Go, humanity spent thousands of years codifying heuristics like ``corners first, then edges, the center last.'' The first generation of AlphaGo still used human game records. AlphaGo Zero did not even need game records---starting from scratch, it surpassed its predecessor AlphaGo in a matter of days.

\vspace{1em}

\textbf{Where is the bitterness?}

The bitterness is this: human pride was crushed.

\vspace{1em}

We believed our intuition, experience, and wisdom were irreplaceable. It turned out that given sufficient compute, machines can rediscover from scratch the patterns it took us millennia to formulate---and even discover patterns we never imagined.

\subsection{Complexity Is Overfitting; Simplicity Is Blank Space}

Why does simplicity always win?

Because complexity is a hidden debt.

Every additional rule is one more place where things can go wrong; every additional module is one more interface to maintain; every additional layer of abstraction is one more assumption that may become obsolete.

Complex systems look powerful, but they typically bring three fatal weaknesses:

\begin{itemize}
\item Fragility---the slightest change in environment, and meticulously tuned parameters fail.

\item Path dependence---once you walk the path of complexity, it is very hard to turn back.

\item Non-scalability---human expert knowledge is finite, while compute can be stacked endlessly.
\end{itemize}

\vspace{1em}

And simple systems?

Simplicity means \textbf{generality}. One algorithm can be applied across images, language, and many other domains.

Simplicity means \textbf{antifragility}. Fewer rules make a system more robust, with fewer points of failure.

Simplicity means \textbf{scalability}. The core mechanism of the Transformer is not complicated, yet by stacking identical structures, it scales from millions of parameters to trillions.

\vspace{1em}

\textbf{Complexity is overfitting to the present. Simplicity is blank space left for the future.}

\subsection{The Price of Simplicity: Waiting for Its Era}

Simplicity is not without cost. The price of simplicity is \textbf{time}.

Waiting for Moore's Law to make compute cheap enough. Waiting for infrastructure to mature.

\vspace{1em}

Many seemingly new technologies are really old methods that waited for their era to arrive.

The idea of neural networks can be traced back to the 1940s, but it was not until around 2012 that they truly exploded, powered by big data, parallel hardware, and public benchmarks. The key ideas behind backpropagation appeared much earlier, yet they too had to wait for engineering conditions to ripen before becoming mainstream.
The Transformer was published in 2017, but it took GPT-3 for the world to realize its power.

\vspace{1em}

If we hold a principle that is both simple and correct, then we can afford to be patient.

\newpage
\section{The Second Thread: Flow Defeats Hoarding}

% \begin{myquote}
% 上善若水。水善利万物而不争，处众人之所恶，故几于道。

% \hfill ——《道德经》
% \end{myquote}

Earlier we discussed Grandet's dilemma: the more valuable data becomes, the less one dares to let it flow.

This is the hallmark of a hoarding mindset.

The hoarding mindset says: ``What I clutch in my hands is mine.''

The flow mindset says: ``What flows through me is mine.''

\vspace{1em}

History has proven it again and again: \textbf{hoarders always lose to those who let things flow.}

Record labels guarded their copyrights---Spotify's streaming upended them.

Encyclopedias guarded their expert-curated entries---Wikipedia's crowdsourcing drowned them.

Taxi companies guarded their medallions---Uber's platform dissolved them.

Banks guarded their branch networks---Alipay's QR codes eroded them.

\vspace{1em}

This is not a moral judgment. It is a law of physics.

In thermodynamics, the total entropy of an isolated system tends to increase.

Open systems like living organisms maintain low-entropy structures locally precisely by continuously exchanging energy and matter with their surroundings.


\begin{importantbox}
Walls can block invaders, but they can also trap you inside.
\end{importantbox}

\subsection{The Real Winners Are Not Owners but Connectors}

In the Fifth Dimension, the real winners are not owners but connectors.

Google does not produce content, but it connects all content.

Uber does not own cars, but it connects all cars.

Airbnb does not own properties, but it connects all properties.

OpenAI does not own human knowledge, but it connects all human knowledge.

\vspace{1em}

What do these companies have in common?

They are \textbf{hubs of flow}, not \textbf{warehouses of stock}.

Their value lies not in what they lock up, but in what they let flow through.

\vspace{1em}

\newpage
\textbf{The} conclusion that follows may be counterintuitive:

\textbf{In a world of flow, openness is safer than closure.}

Closure means relying only on your own resources; once they run out, you are finished.

Openness means becoming a node in the network; as long as the network exists, there is an endless stream of replenishment.

\subsection{Don't Sit on a Gold Mine---Become a Smelter}

So how do you build a moat in a world where everything flows?

The answer is: \textbf{don't sit on a gold mine---become a smelter.}

A gold mine is solid---once mined out, it is depleted, and anyone else can mine it too.

A smelter is fluid---ore flows in continuously, refined steel flows out continuously.

\vspace{1em}

If we can smelt others' raw materials into higher-value products, the raw materials will come to us of their own accord. Not because we monopolized anything, but because we \textbf{created value}.

\vspace{1em}

This is why the most successful AI companies are often not the ones hoarding data, but the ones that set data in motion and refine value from the flow.

\begin{importantbox}
Don't ask what I own. Ask what is willing to flow through me.
\end{importantbox}

\newpage
\section{The Third Thread: Self-Evolution Defeats Stasis}

If we had to pick the most powerful algorithm in Earth's history, the answer would be neither deep learning nor quantum computing, but an ancient program that has been running for billions of years---\textbf{self-evolution.}

Evolution needs no designer. It requires only three elements:

\textbf{Mutation}---a mechanism for generating variation.

\textbf{Selection}---a mechanism for keeping the good and discarding the bad.

\textbf{Inheritance}---a mechanism for passing good traits forward.

\vspace{1em}

It was these three elements that, from a pot of primordial soup, gave rise to language, thought, and civilization.

No blueprint. No plan. No intelligent design.

Only \textbf{mutation, selection, inheritance}, cycling for hundreds of millions of years.

This is the essence of the third thread: systems capable of self-evolution will ultimately surpass static designs.

Not because evolution is smarter,
but because evolution operates on an immense timescale, can try a vast number of possibilities, and will never voluntarily stop.

\subsection{Evolution Is Everywhere}

Evolution does not occur only in the biological world. Systems that possess the three elements of mutation, selection, and inheritance tend to exhibit evolutionary characteristics.

\vspace{1em}

\textbf{The scientific method} is a form of evolution.

Hypotheses are mutations---scientists propose various theories.

Experiments are selection---theories that contradict reality are eliminated.

Publication is inheritance---validated theories enter textbooks and become the starting point for the next generation of research.

This is why science can progress continuously.

It does not depend on individual genius but on a \textbf{self-correcting} collective mechanism.

\vspace{1em}

\textbf{Language and culture} are also forms of evolution.

New words, new memes, new ideas constantly emerge---this is mutation.

Some catch on; some die out---this is selection.

Those that catch on are imitated, spread, and recorded---this is inheritance.

\vspace{1em}


\textbf{Open-source communities} likewise exhibit evolutionary characteristics.

Code forking is mutation---anyone can modify the code.

Usage is selection---useful versions attract more users.

Merging is inheritance---verified improvements are merged back into the main branch.

This is also why Linux has long held a dominant position in the server domain.

Not because any single engineer is smarter, but because the entire community is \textbf{evolving in parallel}.




\vspace{1em}

\textbf{The market economy} can be seen as a form of evolution.

Businesses are mutations---countless entrepreneurs try different business models.

Prices are selection---consumers vote with their wallets, eliminating inefficient firms.

Capital is inheritance---successful models attract more investment and continue to expand.

This is the deeper reason why market economies outperform planned economies---a planned economy is static design; a market economy is continuous evolution.

\vspace{1em}

\textbf{The immune system} is likewise a form of evolution.

Cells randomly generate antibodies---this is mutation.

Antibodies that recognize pathogens are activated---this is selection.

Memory cells preserve effective antibody templates---this is inheritance.

This is why the immune system can respond to many pathogens it has never encountered before.

It does not need to know what the enemy is. It only needs to \textbf{keep evolving}.

\begin{importantbox}
A system that is open enough and enduring enough, under the pressure of survival, will spontaneously move toward evolution to adapt to environmental change.
\end{importantbox}

\newpage
\subsection{The Goal of Evolution: Leaving More That Endures}

What exactly is evolution optimizing for?

Not strength, not appearance, not intelligence---but \textbf{leaving behind more copies or descendants that can persist}.
From the gene-centered perspective, life appears to be doing one thing: making more copies of itself to stave off annihilation. This is the point Dawkins's ``selfish gene'' seeks to emphasize---genes do not care whether their host is happy; they ``care'' whether they can be copied forward.

\vspace{1em}

A gene that helps its host live longer and reproduce more will spread through the gene pool.

A gene that causes its host to die young or become infertile will vanish from the gene pool.

It is that simple.

\vspace{1em}

But replication itself raises a problem: the environment is not uniform.

The Arctic is not the tropics. The ocean is not the land. Today is not tomorrow.

If all copies were identical, a single environmental upheaval would wipe them all out.

\vspace{1em}

So evolution invented mutation.

Not all copies are the same; each generation carries some random changes.

Most changes are harmful and get eliminated.

A few changes are beneficial and spread in the new environment.

It is like venture capital---trading a multitude of failures for a handful of successes.

\vspace{1em}

Seen this way, life is like software: constantly iterating, constantly refactoring.

Version 1.0 adapted to a certain environment and replicated in large numbers.

The environment changed; version 1.0 began to fail.

Certain mutant 2.0 versions were better suited to the new environment and began to spread.

Version 2.0 replaced version 1.0 and became the new mainstream.

Then the environment changed again, and version 3.0 entered the stage\ldots

\vspace{1em}

Every version upgrade is a process of \textbf{pruning redundancy}.

Features that were once useful but are now useless get trimmed away by evolution.

Because maintaining that redundancy costs energy, and energy is scarce.

\textbf{Evolution is a ruthless code reviewer.}

From this we should recognize an important corollary: nothing can live forever.

Eternity belongs only to systems that continuously refactor themselves.

\vspace{1em}

Individual bacteria have extraordinarily short lifespans, yet the species has existed for billions of years.

Their iteration speed is extreme.

Under favorable conditions, some bacteria can complete a round of division in mere minutes.

No matter what environmental change occurs, some mutant version will survive.

This is why antibiotics lose their effectiveness---

It is not that bacteria learned to resist; it is that resistant mutant versions were \textbf{selected} out.

We wiped out the vast majority of bacteria, but the few with drug resistance kept replicating.

Within hours, their numbers rebounded---and every single one was a resistant version.

\vspace{1em}

Companies are no different.

Kodak invented the digital camera, but Kodak fell---

because it refused to refactor itself. It tried to make its version 1.0 immortal.

\textbf{It was replaced by version 2.0 competitors.}

The sensible survival strategy is: \textbf{before others make you obsolete, proactively reinvent your old version.}

Steve Jobs once said: ``If you don't cannibalize yourself, someone else will.''

This is not just business wisdom. It is the fundamental logic of evolution.

\vspace{1em}

So do not pursue immutability. Do not pursue eternity. Do not pursue a final version.

Today's self should be version 2.0 of yesterday's self.

Tomorrow's self should be version 3.0 of today's self.

% \vspace{1em}

% 这就是进化的终极智慧：

% \begin{importantbox}
% 不死的不是任何一个版本，而是版本更新本身。

% 不朽的不是任何一个形态，而是形态变化的能力。
% \end{importantbox}

\subsection{The Enemies of Evolution}

What kind of system cannot evolve?

\textbf{A system without mutation.} If all individuals are identical, there is no room for selection. Inbreeding leads to species decline precisely because of the loss of genetic diversity.

\textbf{A system without selection pressure.} If good and bad are indistinguishable, if superior and inferior cannot be told apart, then evolution loses its direction. Enterprises stagnating under an egalitarian ``iron rice bowl'' system is precisely the result of lacking competitive pressure.

\textbf{A closed system.} Without external input, the sources of mutation dry up. This is why \textbf{flow} is a precondition for evolution---closed systems can only spin within a limited space of possibilities.

\textbf{An overly complex system.} If a system is too complex, any mutation may cause it to collapse. This is why simplicity is a precondition for evolution---complex systems are too fragile to withstand trial and error.

\section{LLMs: The Dragon Slayer and the Dragon}

% \begin{myquote}
% He who fights with monsters should look to it that he himself does not become a monster.

% \hfill ——尼采
% \end{myquote}

LLMs defeated the previous generation's complex systems through ``simplicity + scale.''

But as they themselves became colossal entities deeply dependent on human data, they became the new complexity, the new rigidity, the new inability to evolve. The dragon-slaying youth who once triumphed through raw compute, now extremely dependent on human text data, is hitting a wall.

\vspace{1em}

Viewed through the completeness rate we introduced in Chapter Two, this wall is not ``not enough data'' but ``the distribution is not right enough.'' For tasks within the distribution, human text already provides an extremely high completeness rate, which is why LLMs can perform astonishing compression, paraphrase, averaging, and standardized answering. Yet what rewrites an era is often not that high-density plain, but a small number of critical out-of-distribution signals.

\vspace{1em}

The question thus becomes how to identify which data points, amid a sea of i.i.d. data, are the effective out-of-distribution ones. The real hurdle for the next generation of AI is not to squeeze the old distribution even more thoroughly, but to let a small number of critical out-of-distribution data points be prominently expressed---without destroying the capabilities already learned from the old distribution. Whoever solves this contradiction first will be best positioned to carry the baton forward.

\subsection*{Is the Scaling Law Dead?}

My answer is \textbf{no}. This is more like a Moore's Law--style relay race. Looking back at the history of Moore's Law, we see that it was never carried by a single technology. Enormous sums of capital, talent, and successive waves of new methods kept pouring in, giving it long-lasting vitality. What sustains the curve is never a single path, but a continuous handoff of new approaches.

\vspace{1em}

When the dividends of pretraining are exhausted, the torch may pass to reasoning, tool use, environmental feedback, or places we cannot yet name today. What new technologies will appear after that? I do not know, but I believe in the power of self-evolution.



\newpage
\section{Last Man Standing}

\begin{myquote}
And how can we bring little children into this dark world we've created and let them be surrounded by people who've given up?

\vspace{1em}

So even if this is the end of humanity as we know it, let's fight to the very end.

\vspace{1em}


Let's let the children know that there is hope if they get together.

\vspace{1em}

And even if it becomes impossible for anybody, it's better to go on fighting to the end than just to give up and say, OK.

\hfill ——Jane Goodall
\end{myquote}

These words are from an interview with Dr.\ Jane Goodall in her later years.

\vspace{1em}

If there are no eternal winners, if today's dragon slayer becomes tomorrow's dragon---then what should we strive for?

\vspace{1em}

My answer is in the preface of this book: ordinary people who refuse to surrender on the eve of the storm.

Not ``winning,'' but ``not surrendering.'' Not ``standing on the summit,'' but ``still able to stand up.''

When upheaval arrives, still willing to see clearly, still willing to prepare, still willing to bear the cost.

\vspace{1em}

\textbf{Last Man Standing}---the one still standing at the end.

The true meaning of this phrase is not about defeating every opponent.

It is about being knocked down countless times and still choosing to stand up.

% \vspace{1em}

% 加缪讲述过西西弗斯的故事。

% 诸神惩罚西西弗斯，让他把一块巨石推上山顶。

% 每次快到山顶，石头就会滚落，然后他必须重新开始。

% 如此循环，直到永恒。

% 这看起来是最绝望的处境。

% 但加缪说："我们必须想象西西弗斯是幸福的。"

% \vspace{1em}

% \textbf{为什么？}

% 因为在推石上山的过程中，他是\textbf{自由的}。

% 他选择推，选择不屈服。

% 他无法控制石头是否会滚落，但他可以控制自己是否继续推。

% 这个选择本身，就是他的尊严。

\vspace{1em}


At the end of this book, I want to read the four lines from the cover one more time.

\textbf{Reaching Distant Peaks.}
This is not merely a triumph of connectivity technology; it also reminds us to perceive a more complete data world. Only when we can see and reach more distant mountains will we stop fixating on immediate gains and losses.

\vspace{1em}

\textbf{Discerning the Infinitesimal.}
How do the sword and shield of data attack and defend under the gaze of the all-seeing eye? How should data flow? This reminds us to identify those important, subtle signals: where the costs lie, where the deceptions hide.

\vspace{1em}

\textbf{Tracing the Past.}
A machine predicts the next token from the preceding $n$ tokens, forecasting the future from yesterday. The more powerful the personalized long-term memory a machine can possess, the more capable it becomes of self-iteration and evolution.

\vspace{1em}

\textbf{Shaping the Future.} This is not an optimistic slogan. The future will not automatically become a better place. It demands choices, training, sacrifice---it demands that we reshape ourselves again and again.

\vspace{1em}

On the eve of the storm, if we refuse to surrender, then these four lines can serve as our guide to action:
\textbf{Reach farther. Look closer. Remember longer. Act first.}

\vspace{1em}

Technology will change. Tools will change. Even \textbf{the definition of what it means to be human} will change.

The snow-capped mountain is always there.

The challenge always remains.

We keep climbing.

\textbf{This is our choice. This is our dignity. This is---what it means to be human.}

\vspace{2em}
\begin{flushright}
\textit{On the eve of the storm}

\vspace{1em}

{\small Perhaps it is the eve of a storm. Perhaps it is merely a long twilight.\\
But whichever it is, those who are prepared will have no regrets.}
\end{flushright}